\section{Marco Teórico}
\label{sec:marco-teorico}

\subsection{Design Science Research}

Design Science Research (DSR) es un paradigma de investigación en sistemas de información orientado a
crear y evaluar artefactos que resuelven problemas organizacionales identificados
\cite{hevner2004,peffers2007}. A diferencia del paradigma conductual (behavioral science), que busca
explicar o predecir fenómenos, DSR busca construir un artefacto --- en este caso, un sistema de software
--- y evaluarlo empíricamente contra el problema que motivó su creación. Este trabajo instancia
específicamente el modelo de proceso de seis actividades de Peffers et al. \cite{peffers2007}
(Sección~\ref{sec:metodos}).

\subsection{Ingeniería de requisitos}

La ingeniería de requisitos de este proyecto sigue el estándar ISO/IEC/IEEE 29148:2018
\cite{iso29148}, que define el ciclo de vida de descubrimiento, elicitación, análisis, especificación,
verificación y gestión de requisitos. Las historias de usuario se documentan en formato Connextra
(``Como/Quiero/Para''), evaluadas contra el criterio INVEST (Independent, Negotiable, Valuable,
Estimable, Small, Testable), y los casos de uso siguen la plantilla ``fully dressed'' de Cockburn
\cite{cockburn2001}, con niveles de meta explícitos (nube, cometa, mar, pez). La calidad individual y de
conjunto de los requisitos se evalúa contra el \textit{INCOSE Guide for Writing Requirements}
\cite{incose2023}, que define nueve características por requisito y seis del conjunto completo.

\subsection{Medición del software: el enfoque Goal-Question-Metric}

El enfoque Goal-Question-Metric (GQM) de Basili, Caldiera y Rombach \cite{basili1994} estructura la
definición de métricas de software en tres niveles: una meta conceptual (\textit{qué se quiere lograr y
por qué}), preguntas operacionales que refinan esa meta, y métricas cuantitativas que responden cada
pregunta. Este trabajo instancia un esquema GQM completo para la meta de rendimiento del sistema
(Sección~\ref{sec:metodos}).

\subsection{Arquitectura de software: microservicios y el modelo C4}

Aunque este sistema no adopta una arquitectura de microservicios en sentido estricto (es un monolito
modular de tres capas), el vocabulario de servicios desacoplados comunicados por API HTTP proviene de la
definición de Fowler y Lewis \cite{fowler2014}. La documentación de la arquitectura sigue el modelo C4
(Contexto, Contenedores, Componentes), ya presente en \texttt{docs/arquitectura/README.md} desde
entregas anteriores del proyecto.

\subsection{Seguridad de aplicaciones web}

El marco de referencia para la auditoría de seguridad es el OWASP Top~10:2021 \cite{owasp2021}, que
clasifica los diez riesgos de seguridad web más críticos según datos agregados de la industria (control
de acceso roto, fallas criptográficas, inyección, entre otros). La autenticación del sistema usa JSON Web
Token bajo RFC~7519 \cite{jwt2019}, cuyo uso indebido (algoritmos débiles, secretos hardcodeados, ausencia
de expiración) es una fuente documentada de vulnerabilidades reales en aplicaciones web
\cite{lu2018,rezaeinasab2022}.

\subsection{Usabilidad: la System Usability Scale}

La System Usability Scale (SUS), propuesta por Brooke \cite{brooke1996} y validada empíricamente por
Bangor, Kortum y Miller \cite{bangor2008}, es un instrumento estandarizado de 10 ítems en escala Likert
que produce un puntaje de 0 a 100 interpretable contra una escala de grados (A+ a F). Se usa en este
proyecto como el instrumento de medición de usabilidad (Sección~\ref{sec:evaluacion}); a la fecha de este
informe se aplicó a 15 personas, de las cuales 4 tienen fecha de aplicación verificable y las 11
restantes están pendientes de confirmar (ver \texttt{docs/mediciones/sus/SUS-RESULTS.md}).

\subsection{Estadística inferencial no paramétrica en ingeniería de software}

Cuando los datos de rendimiento no siguen una distribución normal (frecuente en mediciones de latencia,
que suelen tener asimetría positiva por colas largas), las pruebas de hipótesis basadas en la
distribución normal (t de Student) no son apropiadas. Arcuri y Briand \cite{arcuri2011} recomiendan el
uso de pruebas no paramétricas como Mann-Whitney U (equivalente a Wilcoxon rank-sum para dos muestras
independientes) junto con una medida de tamaño de efecto, en vez de reportar solo un valor p. Este
trabajo sigue esa recomendación en el análisis de caché fría/caliente (Sección~\ref{sec:evaluacion}).

\subsection{Metodologías empíricas complementarias}

Runeson y Höst \cite{runeson2009} definen las guías de rigor para investigación de estudio de caso en
ingeniería de software, y Wohlin et al. \cite{wohlin2012} las de experimentación controlada --- ambas
metodologías complementan a DSR cuando la evaluación de un artefacto requiere comparar condiciones
(como la comparación de caché fría/caliente de la Sección~\ref{sec:evaluacion}, que sigue la lógica de un
experimento controlado de una sola variable).

\subsection{Estándares empíricos y revisión sistemática de literatura}

La comunidad ACM SIGSOFT propone un conjunto de \textit{estándares empíricos} \cite{ralph2021} que
codifican las expectativas de rigor para distintos tipos de estudio en ingeniería de software (estudio de
caso, benchmarking, investigación de ingeniería, entre otros) --- aplicados en este proyecto en
\texttt{docs/checklists/} y resumidos en el Anexo D (Sección~\ref{anexo:ralph}). Para la revisión de trabajos relacionados
(Sección~\ref{sec:trabajos-relacionados}), se sigue el proceso descrito por Kitchenham y Charters
\cite{kitchenham2007} y se reporta bajo la disciplina PRISMA 2020 \cite{page2021}, complementado con la
técnica de \textit{snowballing} de Wohlin \cite{wohlin2014ease} para expandir la búsqueda más allá de las
cadenas de consulta iniciales.
\newpage
